\chapter{109}

\section{Bettmeralp/VS (CH), im
Oktober}



Mit der Sonne zog sich auch allmählich die Wärme zurück. Der Herbst
drückte den Sommer täglich entschlossener weg.

Ihre abendlichen Gespräche hatten sich über die Tage etabliert. Auch
heute hielt sich die dichte Erzählstimmung. Sie rückten kontinuierlich
näher an die Feuerschale, um sich warm zu halten.

David drehte versunken eine Bierflasche in seinen Händen hin und her. Es
war nicht zu übersehen, wie er an Johs Lippen hing. Sein Gesicht
präsentierte sich immer wieder als offenes Buch. Als sähe man darauf
Gedanken auf und ab hüpfen, bis er sie bändigte und sie zu gesprochenen
Worten wurden.

«Bist du gläubig, Joh,» fragte er zögernd, «ich mein, mit Kirche und
so?»

Mit dieser Frage traf er unmissverständlich ins Schwarze. Louis wurde
hellwach. Joh wirkte konzentriert, schien seine Worte wohlüberlegt
geordnet zu haben, bevor er zu erzählen begann. Von seinen Erfahrungen
als katholischer Messdiener und davon, wie er in jungen Jahren aus der
Kirche ausgetreten war. Rebellisch überzeugt sei er gewesen und dabei
von gleichgesinnten Freunden getragen.

«Meine Eltern verstanden meinen Kirchenaustritt nicht. Mensch, war das
ein Theater. Es verletzte sie masslos, dass ich nun ein
\emph{Ungläubiger} war, ein endgültig Verlorener. Aber es gab zwischen
uns damals weit kompliziertere Unstimmigkeiten. Das Thema verschwand.
Und wurde zum Tabu.»

Louis streckte seinen Rücken durch.

«Erinnerst du dich, Manuela, als du mir in den ersten Tagen nach meiner
Ankunft ein Musikstück empfohlen hast, \emph{«Le Mantra du Coeur»}
hiess es, von einem…?»

«Von «Michel Pépé», ja.»

Ein Scheit fing Feuer.

«Ich war ziemlich durch den Wind.»

«Ja. Ich habe mich immer mal wieder gefragt, was in diesem Moment bei
dir passiert ist, Louis. Das Stück löst in mir nichts als Ruhe und
Frieden aus. Es war unübersehbar, dass es bei dir komplett anders war.»

Manuela suchte im Dunkeln nach Louis’ Gesicht.

Es kostete ihn wenig Mut, seine Kirchengeschichte zu erzählen. Sie
passte gerade einfach in diese Runde, zu diesem Abend, zu dieser
wohlwollenden Stimmung und unter den Schutz der Nacht.

«Ich, ich war noch sehr klein,» begann Louis und alle Blicke richteten
sich auf ihn, «als mein Vater, mich und meine kleine Schwester an einem
Sonntag mit in die Kirche nahm. Mir war, als setze mich dieses
Musikstück zurück auf diese Kirchenbank. Es war so schrecklich, es, hm,
es klingt verrückt, doch diese Musik, dieser leidende Jesus am Kreuz,
ich hatte damals Angst, dass er runterkommt, und ich, ich konnte nicht
weg. Das kam alles hoch, als die Männerstimme in diesem Musikstück
einsetzte.»

Louis fiel es schwer, seine Traurigkeit, die sich mit jedem seiner Worte
unaufhaltsam nach aussen drängte, zu unterdrücken.

Joh erhob sich und legte ein Holzscheit nach. Für eine geraume Zeit
einte sie das unregelmässige Knistern in der Feuerschale.

Bis sich David unerwartet und mit unsicherem Blick an Louis wandte.

«Wo war deine Mutter?» fragte er leise. Es klang, als zweifle er, ob die
Frage angebracht war. Doch Louis kam sie recht. Sie nahm der Situation
die Schwere und half ihm, sich zu fangen.

«Meine Mutter machte sich nichts aus Kirchenbesuchen,» fuhr er in Davids
Richtung fort, «sie kam nie mit. Sie erklärte meiner Schwester und mir,
dort würden Geschichten erzählt, die sich die Menschen ausgedacht
hätten. Und dass es verschiedene Geschichten gäbe. Verschiedene
Religionen. Und auch, dass sie an etwas Grosses glaube, an das Gute, das
jeder Mensch in sich selbst finden könne. Dazu brauche es keine Kirche.
Ihr müsst wissen, ehm, meine Eltern stritten schon damals fast
ununterbrochen. Ich glaube, die waren sich unterdessen nur noch darin
einig, zwei gemeinsame Kinder zu haben.»

Louis schaute unsicher in die Runde. Es berührte ihn, wie aufmerksam sie ihm
zuhörten und offenbar auf eine Fortsetzung warteten.

«Mein Vater hatte es eines Tages entschieden. Er wollte uns Kindern das
Christentum nahe bringen. Vor allem mir. Ich sei alt genug. Er sprach
vom einzig wahren Glauben. Er hatte alte Kindheitserinnerungen
hervorgekramt. Von damals in Spanien. Wir sollten getauft werden.
Endlich, mit allem Drum und Dran. Ich sollte ein Diener des Pfarrers
werden. Ich würde bedeutsam aussehen. In einer heiligen Robe, hat er
gesagt. Mir gefiel das. Meiner Mutter nicht. Ich habe mich später
gewundert, dass sie ihn machen liess,» Louis schüttelte den Kopf, «auch
die Geschichte von der Entstehung der Welt hat mir mein Vater dargelegt,
blumig und überzeugt sprach er von einem Adam und seiner Eva. Den ersten
Menschen auf der Welt. Ich habe begeistert zugehört und gestaunt.
Und später habe ich Maman vehement widersprochen, bis mir die Tränen
kamen, als sie mir die wissenschaftliche Version der
Entstehungsgeschichte des Menschen erklärte. Es, es war wie eine
Entzauberung. Ich weiss noch genau, wie wütend ich auf sie war.
Schliesslich blieb es bei einigen wenigen Kirchenbesuchen. Vaters Idee
war plötzlich verschwunden. Jesus weg.»

Louis atmete auf.

«Nun ja, wen wunderts, zu einer Taufe ist es nicht gekommen.»

Er stöhnte.

«Und so ging es eigentlich mit all seinen Ideen, du hast es neulich gut
zusammengefasst, Joh, gross inszeniert, nichts dahinter und plötzlich
weg. Dann kam das Nächste. Ich denk jetzt gerade, es fühlte sich an, als
bliese er mich als seinen Sohn zu einem kleinen König auf und liess mir
auch wieder hemmungslos die Luft raus.»

Louis war über seine scharfen Worte selbst am meisten erstaunt. Einzig
Joh wusste um sein desolates Verhältnis zu seinem Vater. Und jetzt,
jetzt redete er vor weiteren Aussenstehenden über ihn. Ausführlich und
abfällig. Er pausierte und schluckte und fühlte sich plötzlich wie ein
Verräter. Durfte er in dieser Weise über ihn herziehen? Hatte er es doch
über Jahre geschafft, seinen Vater zu vergessen, abzuhaken, waren nun
die Geschichten wie angestossene Steine aus ihm herausgerollt.

Ein unheimliches Gefühl beschlich ihn. Was sagten seine Äusserungen über
ihn selbst aus? Würden sie ihn als einfältigen Sohn dieses Mannes sehen?
Der Gedanke beschämte ihn. Wer konnte er anderes sein, als dieser Sohn?
Was war nach seiner Abkehr von Vater überhaupt noch von Louis
übriggeblieben?

Und was dachten sie über Maman?

\sectionbreak

Louis hatte nicht bemerkt, dass Manuela aufgestanden war. Sie stand 
plötzlich neben Joh. Die beiden schauten sich wortlos an und nickten
einander kaum sichtbar zu. Manuela verschwand unauffällig Richtung
Haus. Einmal mehr beobachtete Louis, wie sie sich, offensichtlich in
Abstimmung mit Joh, unerwartet zurückzog.

Louis blickte an sich hinunter, dann irritiert hoch. Er hatte die Worte
ausgesprochen und nun lagen sie da. Es war ihm fremd, dass keiner seine
Erzählung in Frage stellte. Dass weder Joh noch David etwas sagte. Er
sollte versuchen, einen Abschluss zu finden.

«Das Drama wurde durch meine Mutter zusätzlich angeheizt, weil sie
endgültig die Trennung wollte. Heute glaube ich, Vater suchte aus Wut
auf sie versessen nach Projekten oder Hobbies für mich, die meiner
Mutter missfielen. Vielleicht wollte er sie bewusst verletzen, sie
dominieren, ihr zeigen, dass er mich auf seiner Seite hatte,» Louis
drehte sein Gesicht in Johs Richtung, «ja, es gelang ihm immer wieder,
mich mit seiner unbändigen Euphorie zu gewinnen und gegen sie
aufzubringen. Der hatte so viel Energie, wenn ich daran denke, wird mir
schlecht.»

Mit jeder Aussage war die nächste ganz von alleine aus Louis
hinausgeschossen gekommen. So schnell, dass ihm schliesslich keine Zeit
blieb, sich darüber zu wundern.

Joh hatte sich unterdessen in Louis’ Richtung vorgebeugt. Ihre Blicke
trafen sich und Louis glaubte, in seinen Augen zu versinken. Johs
Präsenz und sein gelegentlich beipflichtendes Nicken gaben Louis
Sicherheit. Wieder fühlte es sich an, als würde er als der Kleine von
damals vor ihm stehen, der er in jungen Jahren war und den er über die
lange Zeit mit sich herumgeschleppt hatte. Dieses kleine Kind, als eine
undefinierbare Last, bis heute an ihm festgekrallt, als viel zu schweres
Gepäck, das nicht abzuschütteln war.

«Dem ging’s gar nicht um mich,» er schluckte und suchte nach Worten, «er
war ein, ein Pseudoengel, das trifft’s, aber als Kind verstehst du das
einfach nicht. Er war mein Vater, einen andern hat man ja nicht. Ich
dachte lange, meine \emph{Mutter} sei verkehrt. Das ist das Traurigste
an allem, eigentlich war er ein verdammter Esel.»

Louis blickte um sich. Manuela war nicht mehr zurückgekommen.
Er schüttelte ungläubig den Kopf. Verrückt, was er den beiden Männern 
alles von sich erzählt hatte. Er fühlte sich seltsam leichter und konnte 
ein verhaltenes Lachen nicht mehr zurückhalten. Auf Davids Gesicht nahm er ein Schmunzeln wahr,   
und Joh schob seinen Stuhl näher an ihn heran. 
Er legte Louis den Arm über die Schulter und drückte ihn an sich.

«Das war schön, Louis, du hast Mut,» jetzt wandten sie sich einander
ganz zu, und Louis hätte beinahe geweint, «es ist schwer, den eigenen
Vater in seiner Unvollkommenheit zu sehen. Du bist nicht der Einzige,
der da durch muss. Nicht wahr, David, das kennen wir auch?»

David verrührte beipflichtend die Hände und wirkte, als sei er stolz
darauf, durch Johs Aussage in diese Runde zu gehören. Dann sah David
wieder nachdenklich aus.

«Pseudoengel, wie meinst du das?» fragte er und runzelte dabei die
Stirn.

Louis rückte sich zurecht und spürte, wie Druck in seinen Kopf stieg.
Die Erklärung kam ihm nur schwer über die Lippen.

«Er war für mich als Kleiner eine Art Schutzengel, dem ich vertraute,
den ich bewunderte, aber eines Tages war nichts mehr übrig davon und
plötzlich lag er als leere Hülle vor mir.»

Johs \emph{«Ja.»} folgte einem kräftigen Druck um Louis’ Schultern. Es
hörte sich beipflichtend an. Joh reckte seinen Kopf etwas zur Seite und
suchte Louis’ Blick.

«Wenn du magst, Louis, können wir zwei uns gerne ein anderes Mal weiter
über unsere Väter unterhalten. Es ist spät geworden, morgen steht viel
an,» Joh blickte erst zu Louis, dann zu David, «als Abschluss biete ich
euch an, uns gemeinsam den Song einer Schweizer Rockband anzuhören. Du
kennst die Gruppe ja, Louis,» und zu David gewandt, «und du kennst
mindestens den Gotthardtunnel.»

Joh zwinkerte David zu.

«Es geht um das Vaterdrama des Sohnes. Steve Lee war auf \emph{seine}
Art davon betroffen. Der Song beschäftigt sich mit der Beziehung von
Vätern und ihren Söhnen. Vor Jahren tröstete er mich über
schwierige Zeiten. Ich litt bis weit ins Erwachsenenalter hinein, weil
ich nicht wusste, was ich tun sollte, um meinem Vater zu gefallen. Ich
hatte das Dauergefühl, ihm nicht zu genügen. Viel, viel später verstand
ich, dass auch er gekämpft hatte und sich vorwarf, als Vater versagt zu
haben. Leider gab es selbst zu dieser Zeit noch kein Happy-End, aber es
rückt näher, auch wenn er nicht mehr da ist. Jedenfalls, ehm, ich wünsch
es mir.»

Mit einem Lächeln ging Joh versunken auf den Vorplatz hinauf und kam mit
dem Lautsprecher zurück. Louis war der Song-Titel unbekannt. Dennoch
freute er sich darauf und auch darüber, wie gut er sich gerade fühlte.
Seine Einschätzung bewahrheitete sich immer wieder neu. Um so mehr, seit
er Ciro endlich hatte verabschieden können. Joh war ein guter Typ. In
seiner Nähe fühlte sich Louis richtig.

Joh setzte sich wieder, zog sein Handy aus der Brusttasche seines Hemds
und rückte seine Brille nach oben. Dann suchte er auf Spotify nach dem
Song und wählte den Boomer an.

Die Musik legte sich gleich einer luziden Schutzhülle über sie. Bald
hingen ihre drei Blicke an den kleinen Flammen in der Feuerschale. Und
diesmal liess Louis die Tränen laufen. Sie kühlten seine überhitzten
Wangen.

\qrblock{Gotthard \emph{«Father is that enough?»}}{media/QR-Codes/053_Father_is_that_enough_Gotthard_Spotify.png}

Die Klänge sanken weit über der Wiese ab und lösten sich im Dunkel auf.

Nachdenklich stand Joh auf. Fast in Zeitlupe goss er aus einem Kessel
Wasser in die Feuerschale. Rauch stieg hoch. Die jungen Männer schauten
ihm zu.

Dann drehte sich Joh zu ihnen um.

«Ich schäme mich heute noch, obwohl, Scham ist vielleicht der falsche
Ausdruck. Meine Kinder, vor allem Luca, mein Sohn, musste mich immer
wieder als Abklatsch meines Vaters aushalten,» und jetzt hatte seine
Stimme etwas Brüchiges, «ich konnte es einfach nicht besser. Und ich
kann euch sagen, selbst wenn ich es schon länger will, ich habe einen
Heidenrespekt davor, mit ihm über diese früheren Zeiten zu sprechen.»

David stand auf. Er ging entschlossen auf Joh zu, stellte sich vor ihn
und richtete sich gerade auf.

«Zu mir würdest du jetzt sagen: Was könnte denn Schlimmes passieren,
oder?»

Joh schaute David verblüfft an, stellte den Kessel ab und schloss ihn in
die Arme.
